% Cover letter using letter.cls
\documentclass{letter}
\topmargin=-1in
\textheight=8in
\oddsidemargin=0pt
\textwidth=6.5in

\begin{document}

\signature{N. Ohayon Rozanes}
\longindentation=0pt
\let\raggedleft\raggedright

\begin{letter}{}

  \begin{center}
    {\large\bf N. Ohayon Rozanes} \\
    {611 Granbury Dr. \\ Allen, Texas, 75013 \\ 469-560-0141}
  \end{center} \vfill

  \opening{Dear Hiring Manager:}

  \noindent I am writing to apply for the AI/ML Intern role at Rocket
  Lawyer for Summer 2026. I am pursuing a double degree in Data Science
  and Mathematics at the University of Texas at Dallas, with most of my
  applied work concentrated in machine learning, Python data pipelines,
  and building software systems around models. Rocket Copilot is the
  kind of product I find genuinely interesting to work on: a deployed
  AI system with real users and a concrete domain, where the gap between
  a model performing well in evaluation and performing well in practice
  is a real engineering problem.

  \noindent The majority of my ML experience has come from two directions.
  The first is research: in my current bioinformatics role I design and
  evaluate unsupervised models over high-frequency behavioral data,
  which has meant a lot of iteration on feature engineering, model
  architecture, and experiment infrastructure. I built the full pipeline
  from raw sensor data through preprocessing, model evaluation, and
  result tracking, and a large part of the job has been learning to
  distinguish between a model that looks correct and one that actually
  is. The second direction is independent project work. I built TradeX,
  a quantitative trading system that trained and backtested PyTorch
  models on market data, which required building real training and
  evaluation pipelines rather than just fitting a model to a dataset.
  Both have given me a good sense of where ML systems actually fail,
  which I think is a more useful frame than just knowing how to train
  a model.

  \noindent I am particularly interested in the LLM side of what Rocket
  Lawyer is doing. Generative models are mathematically interesting to
  me, and I have spent time outside of coursework reading about how
  they actually work rather than just how to use them. Working on
  prompting, evaluation, and fine-tuning in the context of a legal
  assistant is a legitimately hard problem: the domain is narrow, the
  failure modes are consequential, and the standard of ``good enough''
  is high. That kind of constraint makes for better engineering, and I
  would rather work on a hard AI problem tied to a real use case than
  an easier one in a more abstract setting.

  \noindent I would welcome the opportunity to contribute to Rocket
  Copilot and to work alongside the team building it.

  \closing{Sincerely yours,}

\end{letter}

\end{document}
